\chapter{Mixed review: five sample-style questions}
This is a fresh practice set modeled on the sample's Q1--Q5. It is not a mock
with a verified current format, timing, or grading scheme. First attempt without
earlier solutions, GHCi, or the reference. Record any assistance afterward.

\task{Question 1. Most general types or ill-typed}
Use the State interface from the reference. All parts are independent. Include
constraints and briefly justify any type error.
\begin{enumerate}[label=(\alph*)]
\item \code{\xs -> [show x | x <- xs, even x]}
\answer{[Answer.]}
\item \code{get >>= \xs -> return (length xs)}
\answer{[Answer.]}
\item \code{modify (map Just)}
\answer{[Answer.]}
\item \code{put True >> return 'k'}
\answer{[Answer.]}
\item \code{(:) <$> get}
\answer{[Answer.]}
\item \code{\f mx -> fmap f mx >>= \y -> return (y,y)}
\answer{[Answer.]}
\end{enumerate}

\clearpage
\task{Question 2. Construct expressions from types}
Give useful total implementations of the described behavior. Use the abstract
State interface. For \code{collectEither}, use explicit recursion, not traversal
helpers. It returns all Right values or the first Left error, left-to-right.
\code{applyAfterRead f} applies \code{f} to current state without changing it.
\code{rotateStore} changes triple state from \code{(x,y,z)} to \code{(z,x,y)}
and returns the old triple. \code{liftThree} performs three actions left-to-right
and combines their results; use do/bind/return, not lifting helpers.
\par\noindent\begin{minipage}{\linewidth}
\begin{codeanswer}
collectEither :: [Either e a] -> Either e [a]
collectEither = undefined  -- YOUR CODE

applyAfterRead :: (s -> a) -> State s a
applyAfterRead = undefined  -- YOUR CODE

rotateStore :: State (a,a,a) (a,a,a)
rotateStore = undefined  -- YOUR CODE

liftThree :: Monad m => (a -> b -> c -> d) -> m a -> m b -> m c -> m d
liftThree = undefined  -- YOUR CODE
\end{codeanswer}
\end{minipage}\par
Can a total function of type \code{Monad m => m a -> a} work for all monads
and all result types? Explain with a counterexample specialization.
\answer{[Answer and reasoning.]}

\clearpage
\task{Question 3. Induced verification conditions}
Use the mathematical function $\operatorname{Squares}(n)=0$ for $n\leq 0$,
and $\operatorname{Squares}(n)=n^2+\operatorname{Squares}(n-1)$ otherwise.
\par\noindent\begin{minipage}{\linewidth}
\begin{lstlisting}
method SumSquares(n:int) returns (t:int)
  requires n > 0
  ensures t == Squares(n) + 4
{
  t := 0;
  var i := n;
  while i > 0 {
    t := t + i*i;
    i := i - 1;
  }
  t := t + 4;
}
\end{lstlisting}
\end{minipage}\par
For each invariant, list initialization, preservation, and exit obligations,
including the suffix after the loop. Derive formulas even for an invalid invariant.
This question asks for partial correctness; termination is outside this set.
\begin{enumerate}[label=(\alph*)]
\item $i\geq0\land t+\operatorname{Squares}(i)=\operatorname{Squares}(n)$.
\answer{Initialization: [ ]. Preservation: [ ]. Exit: [ ].}
\item $i\geq0\land t=\operatorname{Squares}(n-i)$.
\answer{Initialization: [ ]. Preservation: [ ]. Exit: [ ].}
\item $t=\operatorname{Squares}(i)$.
\answer{Initialization: [ ]. Preservation: [ ]. Exit: [ ].}
\end{enumerate}

\clearpage
\task{Question 4. Trace and rewrite a generic program}
\par\noindent\begin{minipage}{\linewidth}
\begin{lstlisting}
pack :: Monad m => [m (a -> b)] -> m [a -> b]
pack [] = return []
pack (u:us) = do
  f <- u
  fs <- pack us
  return (f:fs)
\end{lstlisting}
\end{minipage}\par
Evaluate each expression, leaving functions as \code{id}, \code{not}, etc.
where needed. For (e), both initial and stored functions are Boolean functions.
\begin{enumerate}[label=(\alph*)]
\item \code{pack [Just not]}
\answer{[Answer.]}
\item \code{pack [Just not, Nothing]}
\answer{[Answer.]}
\item \code{pack [[] :: [Bool -> Bool]]}
\answer{[Answer.]}
\item \code{pack [[not,id], [const True, const False]]}
\answer{[Answer in enumeration order.]}
\item \code{runState (pack [get, modify (. not) >> get]) id}
\answer{[Returned functions and final stored function; apply each to True to check.]}
\end{enumerate}
Reimplement \code{pack} using exactly one call to \code{mapM}, without recursion.
\par\noindent\begin{minipage}{\linewidth}
\begin{codeanswer}
packViaMapM :: Monad m => [m (a -> b)] -> m [a -> b]
packViaMapM = undefined  -- YOUR CODE
\end{codeanswer}
\end{minipage}\par

\clearpage
\task{Question 5. Lazy summary and reconstruction}
\par\noindent\begin{minipage}{\linewidth}
\begin{lstlisting}
scatterMin :: [Int] -> Int -> (Int,[Int])
scatterMin [] seed = (seed,[])
scatterMin [x] seed = (x,[seed])
scatterMin (x:xs) seed = (min x rest,seed:out)
  where (rest,out) = scatterMin xs seed
\end{lstlisting}
\end{minipage}\par
Evaluate:
\begin{enumerate}[label=(\alph*)]
\item \code{scatterMin [] 11}
\answer{[Answer.]}
\item \code{scatterMin [8] 11}
\answer{[Answer.]}
\item \code{scatterMin [8,3,6] 11}
\answer{[Answer.]}
\item \code{scatterMin [8,3,6] (-20)}
\answer{[Answer.]}
\end{enumerate}
Implement replacement of every element by the list minimum, returning an empty
list on empty input. Use exactly two calls to \code{scatterMin} for nonempty input.
Do not use explicit recursion, folds, \code{minimum}, \code{map}, \code{replicate},
or other recursive list functions. Then give a one-call version under the same
restrictions using a lazy recursive binding.
\par\noindent\begin{minipage}{\linewidth}
\begin{codeanswer}
toMinimumTwoPass :: [Int] -> [Int]
toMinimumTwoPass = undefined  -- YOUR CODE

toMinimumOneCall :: [Int] -> [Int]
toMinimumOneCall = undefined  -- YOUR CODE
\end{codeanswer}
\end{minipage}\par
\answer{[Explain why the one-call binding is productive on finite nonempty inputs.]}

\task{After your attempt}
Record which questions were independent, hinted, or completed after an
explanation. Choose one type error to revisit with a changed example.
\answer{[Evidence and next review target; no predicted instructor grade.]}
